\documentclass[11pt]{article}
\usepackage{amsmath,amssymb,amsthm,booktabs}
\usepackage[margin=2.5cm]{geometry}
\newtheorem{theorem}{Theorem}\newtheorem{lemma}[theorem]{Lemma}\newtheorem{proposition}[theorem]{Proposition}
\newtheorem{corollary}[theorem]{Corollary}
\theoremstyle{definition}\newtheorem{definition}[theorem]{Definition}\newtheorem{remark}[theorem]{Remark}
\title{Q597 G1: crossing translation (CT) reduced to a wake-shift statement\\
\normalsize PROVE 2026-11-19. CT \emph{not proved}. Proved: CT $\Leftarrow$ WS (a two-sided statement about the tiles the crossing reads).
CT at size 6: CT1 and CT2 hold on 4501/4501 (all candidates; 0 failures, lesion 0).}
\author{Clio}\date{19 November 2026}
\begin{document}\maketitle

\section*{Provenance}
\begin{itemize}
\item Code: \texttt{proofs/code-1119-Q597/} (new; \texttt{ls} first showed no \texttt{code-1119*}). No script shadows a stdlib module. Engine unchanged (\texttt{code-1112c2-Q585/migrate.py}).
\texttt{ct6.py}: CT1/CT2 over all candidates, the lesion, shared track, middle grammar, the $I=\emptyset$ classification, and a $B/B$ control. \texttt{prov.py}: the provenance dictionary (\S\ref{sec:dict}) with its own $B/B$ control.
Both reuse the census and step-log machinery of \texttt{code-1118c2-Q597/bs6.py} unchanged (by \texttt{exec} of its preamble) and \texttt{stepset.py} (validated 11-18 against the engine, 521/521).
\item Regression first: \texttt{ct6.py 5} reproduces the 11-18 numbers exactly. CT1 and CT2 hold on 253/253 ($q$) and 268/268 ($p$), and the shared track holds on 120+136 = 256 of 256 nonempty runs.
\item Setting and notation: \texttt{2026-11-18c2-Q597-BS.tex} \S1 and Def.~CT. No external source is used in this file.
\end{itemize}

\section*{Summary}
\begin{enumerate}
\item \textbf{Step 0.} The 11-18 c2 file is completed (placeholders filled from \texttt{bs6.out}: BS 4501/4501 at size 6, the catalogue grows, G4-by-catalogue dead). The $I=\emptyset$ counter is explained: it is a \emph{reroute with the same endpoints}: $a\,b\,a\leftrightarrow b\,b\,b$ (all 54 at $\le5$; 673 of the 717 stored at size 6), plus a new size-6 shape $a\,b\,b\leftrightarrow b\,b\,a$ (44) whose blocks are not rigid. All of these lie outside BS's hypothesis $I\neq\emptyset$.
\item \textbf{Step 1 (falsifier).} \texttt{ct6.py 6} (33159 states, all candidates). \textbf{CT1 and CT2 hold on 4501/4501} (2447 $q$, 2054 $p$). The lesion passes 0 times. No candidate reads a common tile (4501/4501).
Shared track: 2540/2540 nonempty runs are rigid translates of a block of $x$'s journey, and all have the shape $a\,b\,w\,b$. No falsifier, so the session proceeds to Step 2.
\item \textbf{Step 2.} \emph{Proved} (Theorem~\ref{thm:red}): CT1 at step $i$ follows from the two-sided wake-shift condition WS$_i$. In words, every tile read by a candidate that the crossing did not itself make lies in $M_P$, and its $-\alpha$ translate lies in $Mx_P$, and conversely. The tie-break tile needs no separate hypothesis, because boards are tilings (Lemma~\ref{lem:acr}). CT2 follows from WS$_\mathrm{last}$ in the same way.
\item \textbf{Correction to the brief's route.} Step (ii) of the brief (``on $x$'s wake the two boards are related by the rigid translate of $x$'s journey block'') is \emph{false} as stated.
A wake tile read by the crossing has its $-\alpha$ translate among $x$'s \emph{added} tiles in only about 60\% of cases. Otherwise it is a \emph{common} tile, untouched by $x$ (\S\ref{sec:dict}).
What is true is a rigid shift of the \emph{read region}, ending at the ends of a strip. That is WS, and it is the next gap.
\item \textbf{Controls.} The reduction is a valid implication on every record, so on $B/B$ it is its \emph{hypothesis} that must fail, and it does: WS(M) fails on 722/880 $B/B$ journeys at $\le5$, mostly at a wake tile whose $-\alpha$ translate is absent. The other 158 fail BS by middle length (\S\ref{sec:ctrl}). The lesion $\alpha\mapsto\alpha\pm e_i$ breaks CT1 on every record (0 passes, 521 and 4501 records).
\item \textbf{Retreat G3(b), run.} The other journey is \emph{not} a block swap after inversion (6/521). It is not a two-block swap at any cut in 471/521. G3 stands.
\end{enumerate}

\section{Setting}\label{sec:set}
A \emph{board} $T$ is a finite set of tiles $(\mathrm{kind},\text{vertex set})$ forming a polyhedral subdivision of a fixed planar region (a mosaic, \texttt{2026-11-17c2-MC-pq-exchange.tex}). In particular every edge lies in at most two tiles of $T$.
A step $\sigma=(\mathrm{typ},g,n)$ applies to $T$ if $g\subseteq T$, and then $T\cdot\sigma=(T\setminus g)\cup n$. Translation by $v\in\tfrac12\mathbb Z^4$ acts on tiles, sets of tiles and steps, written $\cdot+v$.
$\mathrm{Cand}(T,R)$, $\mathrm{acr}(T,R)$ and $\mathrm{step}(T,R)$ are as in \texttt{2026-11-18c2-Q597-BS.tex}, Prop.~cand, for a fixed migration direction $d$.

Fix a clean nontrivial $B/C$ record. $M_P$ and $Mx_P$ are the two boards after the common prefix of the clean journey. The $M$-middle is $a\cdot c_1\cdots c_k$. The rhombus is $R_a$ before $a$ and $R_i$ before $c_i$, so $R_1=R_a+\alpha$, and $\gamma=R_{k+1}-R_1$ is the displacement of $C$.
Put $M_1=M_P\cdot a$ and $M_{i+1}=M_i\cdot c_i$. Write $g_a,n_a$ for the tiles removed and added by $a$, and $g_i,n_i$ for those of $c_i$.

\section{The reduction (proved)}\label{sec:red}
\begin{definition}[made by the crossing]
$N_i\subseteq M_i$ is the set of tiles of $M_i$ that some $c_j$ ($j<i$) created, i.e.\ $N_1=\emptyset$ and $N_{j+1}=(N_j\setminus g_j)\cup n_j$. Put $O_i=M_i\setminus N_i$. For any board $X$ to which $c_1-\alpha,\dots,c_{i-1}-\alpha$ apply in turn starting from $Mx_P$, define $N'_i$ in the same way from the translated steps, and $O'_i=Mx_i\setminus N'_i$.
\end{definition}

\begin{lemma}[history translates]\label{lem:hist}
If $Mx_{j+1}=Mx_j\cdot(c_j-\alpha)$ for $j<i$, then $N'_i=N_i-\alpha$ and $O'_i\subseteq Mx_P$. Also $O_i\subseteq M_P\cup n_a$.
\end{lemma}
\begin{proof}
$N'_1=\emptyset=N_1-\alpha$, and $N'_{j+1}=(N'_j\setminus(g_j-\alpha))\cup(n_j-\alpha)=N_{j+1}-\alpha$ because translation is a bijection on tiles.
A tile of $Mx_i$ is either in $Mx_P$ and never removed, or it was created by some $c_j-\alpha$. In the second case take the last creation: the tile was not removed afterwards, since otherwise it would be absent or re-created later, so it lies in $N'_i$. Hence $O'_i\subseteq Mx_P$. The same argument for $M_i$ also allows creation by $a$.
\end{proof}

\begin{definition}[wake shift WS$_i$]\label{def:ws}
Write $\mathrm{Rd}(T,R)=\bigcup\{S\setminus\{R\}:S\in\mathrm{Cand}(T,R)\}$ for the tiles a step reads. WS$_i$ is the conjunction of
\begin{itemize}
\item[(M)] every $u\in\mathrm{Rd}(M_i,R_i)\cap O_i$ lies in $M_P$, and $u-\alpha\in Mx_P$;
\item[(X)] every $u'\in\mathrm{Rd}(Mx_i,R_i-\alpha)\cap O'_i$ has $u'+\alpha\in M_P\setminus g_a$.
\end{itemize}
\end{definition}

\begin{lemma}[tie-break is free]\label{lem:acr}
Let $T,T'$ be boards with $\mathrm{Cand}(T',R+v)=\mathrm{Cand}(T,R)+v$. Suppose $\mathrm{step}(T,R)$ is defined. Then the tie-break keeps translated sets: $\{S\in\mathrm{Cand}(T',R+v):S\cap\mathrm{acr}(T',R+v)\ne\emptyset\}=\{S\in\mathrm{Cand}(T,R):S\cap\mathrm{acr}(T,R)\ne\emptyset\}+v$.
\end{lemma}
\begin{proof}
$\mathrm{acr}(T,R)$ is the set of tiles $\ne R$ of $T$ containing a designated edge $e$ of $R$. Here $e$ is a function of $(R,d)$ alone: the horizontal edge at the unique leftmost vertex in the frame of $d$. So the designated edge of $R+v$ is $e+v$, and the uniqueness condition for $R+v$ is the same as for $R$.
If the uniqueness fails, there is no filter on either side. Otherwise, $T$ is a subdivision, so $\mathrm{acr}(T,R)$ has at most one element.
The filter is applied only when $|\mathrm{Cand}|>1$, and the candidate counts agree. Since $\mathrm{step}(T,R)$ is defined, some $S_0\in\mathrm{Cand}(T,R)$ contains a tile $t\in\mathrm{acr}(T,R)$, so $\mathrm{acr}(T,R)=\{t\}$.
Then $S_0+v\in\mathrm{Cand}(T',R+v)$, so $t+v\in T'$, and $t+v\ne R+v$ contains $e+v$. Hence $\mathrm{acr}(T',R+v)=\{t+v\}$, and $S\ni t\iff S+v\ni t+v$.
\end{proof}

\begin{theorem}[CT1 from WS]\label{thm:red}
Assume WS$_i$ for $i=1,\dots,k$, and that the termination test fails at $(Mx_i,R_i-\alpha)$ whenever it fails at $(M_i,R_i)$ (part of Def.~CT on 11-18). Then for every $i\le k$ the clean journey in the other state performs $c_i-\alpha$ at its $i$-th middle step, and \emph{(CT1)} $\mathrm{Cand}(Mx_i,R_i-\alpha)=\mathrm{Cand}(M_i,R_i)-\alpha$ holds.
If moreover WS$_{\rm last}$ holds, i.e.\ $\mathrm{Rd}(M_P,R_a)+\gamma\subseteq Mx_{k+1}$ and $\mathrm{Rd}(Mx_{k+1},R_a+\gamma)-\gamma\subseteq M_P$, then (CT2) holds, and so (Thm.~ct of 11-18) BS holds for the record.
\end{theorem}
\begin{proof}
Induction on $i$, with the hypothesis that $Mx_i$ is the actual board, i.e.\ $c_j-\alpha$ was executed for $j<i$, and that $y$'s rhombus is at $R_i-\alpha$. For $i=1$ this is the common prefix, which leaves the rhombus at $R_a=R_1-\alpha$.

($\subseteq$) Let $S\in\mathrm{Cand}(M_i,R_i)$ and $u\in S$. If $u=R_i$, then $u-\alpha$ is the rhombus of $Mx_i$. If $u\in N_i$, then $u-\alpha\in N'_i\subseteq Mx_i$ by Lemma~\ref{lem:hist}.
Otherwise $u\in O_i$, and (M) gives $u\in M_P$ and $u-\alpha\in Mx_P$. Suppose some $c_j-\alpha$ with $j<i$ removed $u-\alpha$, and take the largest such $j$. Then $u\in g_j$, so $c_j$ removed $u$ in $M$. Since $u\in M_i$, a later $c_l$ with $j<l<i$ re-created it, so $u\in N_i$, a contradiction. So $u-\alpha\in Mx_i$, and $S-\alpha\subseteq Mx_i$.
The admissibility tests read only the vertex sets of the four tiles of $S$, of $R$, and $d$. Connectivity through shared non-rhombus edges is a property of $S$ alone. All the tests commute with translation (proof of Prop.~cand, 11-18). Hence $S-\alpha\in\mathrm{Cand}(Mx_i,R_i-\alpha)$.

($\supseteq$) Let $S'\in\mathrm{Cand}(Mx_i,R_i-\alpha)$ and $u'\in S'$. The rhombus and $N'_i$ are handled as above. If $u'\in O'_i$, (X) gives $u'+\alpha\in M_P\setminus g_a$, so $a$ does not remove it.
If some $c_j$ ($j<i$, largest) removed $u'+\alpha$, then $c_j-\alpha$ removed $u'$, and since $u'\in Mx_i$ it was re-created, so $u'\in N'_i$, a contradiction. So $u'+\alpha\in M_i$, and $S'+\alpha\in\mathrm{Cand}(M_i,R_i)$ as before.

So CT1 holds at $i$. By Lemma~\ref{lem:acr}, the termination hypothesis and Prop.~cand, $\mathrm{step}(Mx_i,R_i-\alpha)=c_i-\alpha$. This is the next actual step, the rhombus moves to $R_{i+1}-\alpha$, and the induction continues.
CT2 is the same argument with no history: the tiles of a candidate translate tile by tile by WS$_{\rm last}$.
\end{proof}

\begin{remark}[what the theorem is, honestly]\label{rem:honest}
Theorem~\ref{thm:red} is close to an equivalence. If CT1 holds at every step, then every read tile $u\in O_i$ has $u-\alpha\in O'_i\subseteq Mx_P$. So CT1 implies WS$_i$ apart from the clause ``$u\in M_P$'', which excludes tiles made by $a$, and its mirror ``$\notin g_a$''.
The theorem is therefore a \emph{localisation}, not a proof of CT. It removes the history (the crossing's own tiles translate automatically, Lemma~\ref{lem:hist}) and the tie-break (Lemma~\ref{lem:acr}).
What remains is a statement about \emph{the tiles of $M_P$ and $Mx_P$ that the crossing reads}. Since $M_P$ and $Mx_P$ differ exactly by $x$'s journey (after the same prefix), WS is a statement about $x$'s wake at the place where $y$ crosses it, which is the form the brief asked to have named.
\end{remark}

\section{The dictionary: what WS says on the data}\label{sec:dict}
\texttt{prov.py} classifies every non-rhombus tile of every candidate on both sides of every crossing step. In $M_i$ the classes are: made by $a$, made by some $c_j$ ($\in N_i$), \emph{wake} ($\in M_P\setminus Mx_P$), \emph{common} ($\in M_P\cap Mx_P$).
In $Mx_i$ they are: made by $c_j-\alpha$, \emph{xadd} ($\in Mx_P\setminus M_P$), common. The translate of each tile is also classified on the other side, or marked absent.
\begin{center}\begin{tabular}{llrr}\toprule
step & tile in its board $\to$ translate in the other & $q$ & $p$\\\midrule
$C$, $M\to Mx$ ($-\alpha$) & $c$-made $\to$ $c$-made & 864 & 1344\\
& wake $\to$ xadd & 1224 & 1481\\
& wake $\to$ common & 825 & 814\\
& made by $a$, common, or $\to$ absent & 0 & 0\\
$C$, $Mx\to M$ ($+\alpha$) & $c$-made $\to$ $c$-made & 864 & 1344\\
& xadd $\to$ wake & 1224 & 1481\\
& common $\to$ wake & 825 & 814\\
& $\to$ absent & 0 & 0\\
last, $M\to Mx$ ($+\gamma$) & common $\to$ xadd / wake $\to$ xadd & 253 / 506 & 268 / 536\\
last, $Mx\to M$ ($-\gamma$) & xadd $\to$ common / xadd $\to$ wake & 253 / 506 & 268 / 536\\\midrule
candidates per $C$-step & 1 / 2 & 395 / 288 & 317 / 448\\\bottomrule
\end{tabular}\end{center}
Counts are tile incidences, $\le5$ census, 521 records (\texttt{prov5.out}). Of the wake reads, the share whose translate is $x$-added is $1224/2049=60\%$ ($q$) and $1481/2295=65\%$ ($p$).
About 60\% of the read wake tiles were displaced by one of $x$'s type-a steps (column \texttt{x-a-disp} in the output); the rest were not, so by Lemma~flipa (lozenge hypothesis, computed 370/370) they were removed by a type-b step of $x$. That last inference was not computed directly.

\paragraph{Size 6: partial run.} \textbf{This is a fixed subsample, not the census: 8 of 24 chunks of the 10855 size-6 states, 1449 records (815 $q$, 634 $p$), \texttt{prov6.out}.} The run finished 12 minutes after the PROVE session ended, and this paragraph was typeset by WAKE (11-19 c2) from that file. The other 16 chunks were not run.
\begin{center}\begin{tabular}{llrr}\toprule
step & tile in its board $\to$ translate in the other & $q$ & $p$\\\midrule
$C$, $M\to Mx$ ($-\alpha$) & $c$-made $\to$ $c$-made & 2901 & 3471\\
& wake $\to$ xadd (of which \texttt{x-a-disp}) & 4130 (2309) & 3698 (2425)\\
& wake $\to$ common (of which \texttt{x-a-disp}) & 2914 (2099) & 2023 (1389)\\
& made by $a$, or $\to$ absent & 0 & 0\\
$C$, $Mx\to M$ ($+\alpha$) & $c\to c$ / xadd $\to$ wake / common $\to$ wake & 2901 / 4130 / 2914 & 3471 / 3698 / 2023\\
last, $M\to Mx$ ($+\gamma$) & common $\to$ xadd / wake $\to$ xadd & 815 / 1630 & 634 / 1268\\
last, $Mx\to M$ ($-\gamma$) & xadd $\to$ common / xadd $\to$ wake & 815 / 1630 & 634 / 1268\\\midrule
candidates per $C$-step & 1 / 2 & 1381 / 967 & 750 / 1157\\\bottomrule
\end{tabular}\end{center}
The share of wake reads whose translate is $x$-added is $4130/7044=58.6\%$ ($q$) and $3698/5721=64.6\%$ ($p$), matching $\le5$. No translate is absent and no read tile is made by $a$.
$B/B$ control in the same subsample (1284 non-commuting pairs): on $x$-journeys the dictionary holds on 569, all with middles of different lengths (3/2: 85, 4/3: 263, 5/4: 216, 6/5: 5), and fails on 715 (97 read an $a$-made tile; 618 have an absent translate, at $C$-steps 0/1/2: 573/41/4). On $y$-journeys it fails on 1284/1284, always by an absent translate, at $C$-steps 0--8.


\begin{itemize}
\item \textbf{Never absent}, on either side (this is CT1/CT2 at tile level), and \textbf{never made by $a$}, so the clause ``$u\in M_P$'' of (M) is not an artefact.
\item \textbf{(M) refined:} a read tile of $O_i$ is \emph{always a wake tile, never common}. So WS(M) splits into (WS-a) \emph{the crossing reads, outside its own tiles, only $x$'s wake}, and (WS-b) \emph{the $-\alpha$ translate of a read wake tile survives $x$'s journey}.
\item \textbf{(WS-b) is not ``the wake translates''.} The translate $u-\alpha$ is an $x$-added tile in roughly 60\% of the reads, and a \emph{common} tile otherwise. The common case is the end of a strip: when a zone of identical translates $u,u-\alpha,u-2\alpha,\dots$ is shifted by $-\alpha$, the interior tiles are common to both boards. Only the first is removed and only the last is added.
This kills step (ii) of the brief's route as stated. The run-zone lemma (11-18 Lemma~runzone) moves \emph{the zone}, not the individual wake tile, by $-\alpha$, and that is the right unit for WS-b.
\item \textbf{(X) refined:} a read tile of $O'_i$ is xadd or common, and its $+\alpha$ translate is always a wake tile of $M_P$ (consistent with (M)).
\item \textbf{CT2:} the $a$-hexagon in $M_P$ reads one common tile and wake tiles, and all their $+\gamma$ translates are $x$-added tiles of $Mx_P$ not touched by $C-\alpha$.
\end{itemize}

\section{Controls}\label{sec:ctrl}
\textbf{Lesion.} $\alpha\mapsto\alpha\pm e_i$ (eight neighbours) in CT1: passes on 0 of 8 for every record, 521/521 at $\le5$ and 4501/4501 at size 6. Since the lesioned recipe still uses the correct steps $c_i-\alpha$ for the board history, the failure is in the candidate comparison itself: the read tiles have their translates at $-\alpha$ only.

\textbf{$B/B$.} Apply the same recipe ($a$ = first middle step, $C$ = the rest, $\alpha$ = displacement of $a$) to both journeys of every non-commuting $B/B$ pair. \textbf{Theorem~\ref{thm:red} is a valid implication there too, so the hypothesis WS must fail, and it does.}
At $\le5$ (\texttt{prov.py}), (M) fails on 722 of 880 journeys. In 687 a read wake tile has its $-\alpha$ translate \emph{absent} (WS-b fails), at $C$-step 0--6. In 35 the step reads a tile made by $a$ (the clause $u\in M_P$ fails).
On the remaining 158 ($x$-journeys only), (M) holds on every compared step but the two middles have \emph{different lengths}. So BS fails there at the level of the journey's shape, through the termination clause of Theorem~\ref{thm:red}, and not through WS(M). (The mirror (X) was not computed on these 158 at $\le5$.)
At size 6 (\texttt{ct6.py}, all candidates): CT1 fails on 6414 of 8158 $B/B$ journeys, and the 1744 where it holds on the compared steps all have middles of different lengths.
\emph{The step of the argument that fails on $B/B$:} WS-b (the wake does not shift rigidly by $-\alpha$ under $y$), or, when it survives, the termination clause. Any proof of WS must therefore use that $x$ and $y$ go to \emph{different} nests.

\section{Where the proof stops (the next gap)}\label{sec:gap}
\begin{description}
\item[G1 = WS.] For a clean nontrivial $B/C$ record, at every step of the crossing: (WS-a) every candidate tile of $y$ in $M_i$ outside $y$'s own crossing tiles lies in $x$'s wake $M_P\setminus Mx_P$; (WS-b) for each such tile $u$, $u-\alpha\in Mx_P$; and the mirror (X).
Proving this needs the geometry of $x$'s swept region. The wake is the union of the hexagons of $x$'s steps, so (WS-a) says that \emph{$y$'s candidate hexagons lie inside $x$'s swept strip}. (WS-b) says that the strip, near the crossing, is shifted rigidly by $-\alpha$.
The available tools are Lemma~flipa and Lemma~runzone (11-18): one type-a run shifts each zone of $x$'s rhombus by a single vector. The shared-track fact (data, $\le5$: 256/256, size 6: 2540/2540) says $y$'s run $w$ is a translate of a block of $x$'s journey, so $y$ travels along $x$'s strip.
\emph{What is missing} is a lemma of the form: if two rhombi are at translates $R$ and $R+\delta$ of each other inside one zone strip, with $\delta$ parallel to the strip, then their candidate hexagons are translates. This is exactly CT1 on the run $w$. The two $b$-turns would then be a finite local check, one per pair of rhombus shapes, independent of size.
That lemma is the uniform statement the catalogue growth demands. It is not proved here.
\item[G3, test (b) (retreat; run, $\le5$).] Is the \emph{other} (non-clean) journey a block swap after inversion? Inverting a journey reverses its steps and exchanges $g\leftrightarrow n$, so its middles are those of the forward journey, read backwards.
BS for the inverse is therefore a two-block swap of the forward middles at the cut $s=L-1$. \texttt{g3b.py} tests every cut $s$: $m'=(m_{[s:]}+v_1)\cdot(m_{[:s]}+v_2)$ with rigid blocks.
\emph{Answer: no.} The cut $s=L-1$ (inversion) works on 6/253 ($q$) and 0/268 ($p$). Some cut works on 50/521 (the other 44 at $s=3$), and \textbf{no cut} works on 471/521. Control: the clean journey has the forward cut $s=1$ on 521/521.
The other middles always have equal length on both sides, and their length histogram is the clean one shifted by $+1$ (a distributional match, not checked record by record). G3 stands. The other journey needs a different description, not a swap of two rigid blocks.
\end{description}

\section{Status}
\begin{itemize}
\item \textbf{Proved:} Lemma~\ref{lem:hist}, Lemma~\ref{lem:acr} (tie-break needs no hypothesis), Theorem~\ref{thm:red} (CT1/CT2 and hence BS from WS). Remark~\ref{rem:honest}: this is a localisation, essentially an equivalence.
\item \textbf{Computed:} CT1, CT2 on all candidates: 521/521 ($\le5$), 4501/4501 (size 6); lesion 0 passes at both; no common tile read: 521/521, 4501/4501; shared track 256/256, 2540/2540; dictionary (\S\ref{sec:dict}) at $\le5$ with no absent translate and no $a$-made read; $B/B$ control as in \S\ref{sec:ctrl}; G3(b) 471/521 no cut.
\item \textbf{Refuted:} the brief's step (ii) as stated (wake tiles translate onto $x$-added tiles); about 40\% of read wake tiles translate onto common tiles.
\item \textbf{Not proved:} CT, WS, BS. Registry: CT node \emph{computed}.
\end{itemize}


\end{document}
